% Chapter 1: Research Description

\chapter{Research Description}
\noindent{This study addresses urban congestion and long commute times in Philippine cities by proposing a Lamarckian memetic algorithm to redesign fragmented jeepney route systems. Utilizing Iligan City as the primary computational environment, the model integrates a parallelized evolutionary algorithm  with an agent-based simulation calibrated with empirical commuter behavior, including walking thresholds and transfer aversion. This study does not consider one-way constraints; instead, the optimization explicitly proposes a complete reconstruction of route systems within the limits of driveable roads. To accelerate convergence, the framework transitions beyond standard Darwinian \ac{GA} by embedding an \ac{ACO} pheromone matrix as an epigenetic memory structure. This Lamarckian approach allows candidate networks to inherit learned passenger demand landscapes, actively pruning tortuous routes and pulling coverage toward high-demand corridors via targeted local search operators. The performance of each candidate network is evaluated in parallel, utilizing Total User Cost (incorporating travel, waiting, and transfer penalties, subject to capacity limits and fleet equity regularizers) as the definitive fitness metric. Validated through systematic benchmarking and theoretical grounding in transit network design, this work establishes the feasibility of combining data-driven travel models with scalable, demand-aware metaheuristics for the modernization of informal public transport.}

% Chapter 1 Section 1.1: Background of the Study

\section{Background of the Study}
\noindent{Jeepneys are the backbone of Philippine public transportation, providing affordable and accessible mobility for millions of commuters daily \citep{mateo-babiano2020}. Despite their ubiquity, jeepney systems suffer from overlapping routes, suboptimal stop placement, and lack of coordination, leading to inefficient service and extended passenger travel times \citep{jica2022, upncts2021b}. The Baguio City jeepney study (2017) documented load factors averaging 0.84, meaning vehicles were operating at 84\% of their seating capacity and utilization ratios of 0.95 in which it indicates routes were active 95\% of the available service time, revealing high operational intensity but limited optimization of route structures. Current modernization policies emphasize vehicle standards and fare structures but lack integrated, data-driven route optimization tools \citep{mateo-babiano2020, jica2022}.

The jeepney route optimization problem addressed in this study is formally situated within the \ac{TNDP}, a well-established and rigorously defined problem class in transportation research. The \ac{TNDP} involves determining an optimal set of transit routes over a city's street network, with the primary inputs being the road graph and \ac{OD} demand between city zones, while constraints govern demand coverage, service level requirements, and resource availability \citep{kepaptsoglou2009, owais2025}. The problem is inherently combinatorial and NP-hard, meaning the solution space grows exponentially with network size such that exact optimal solutions are computationally intractable for real-world urban instances (Kepaptsoglou and Karlaftis, 2009; Wens and Vansteenwegen, 2025). This study adopts the \ac{TRNDP} subproblem of the \ac{TNDP}, wherein the decision variables are the spatial configurations of jeepney routes over Iligan City's driveable road network. Unlike conventional public transit networks that operate on rigid schedules and designated stops, paratransit systems in developing regions are decentralized and uncoordinated, with passengers boarding and alighting dynamically along a route. Applying the TNDP framework to such informal networks requires modeling individual traveler agents who navigate localized capacity constraints and significant transfer penalties. The scope of this study is deliberately bounded: dynamic demand, real-time traffic feedback, and financial cost objectives are excluded; operational parameters such as vehicle capacity and overall fleet size are fixed, while fleet allocation across active routes is determined dynamically via Mohring's square-root rule based on route demand. The NP-hard nature of the \ac{TRNDP} formally motivates the adoption of metaheuristic approaches over exact methods, and specifically justifies the hybrid Memetic \ac{GA} with \ac{ACO}-inspired Lamarckian learning developed in this study as the optimization engine.

\ac{ABM} provides a powerful framework for simulating complex urban mobility systems where individual passenger and vehicle behaviors drive emergent system-level outcomes \citep{mehdizadeh2022, li2025}. Recent advances in \ac{ABM}-\ac{GIS} integration enable realistic spatial grounding of agent decision-making, incorporating \ac{OSM} data, land-use patterns, and behavioral parameters \citep{wozniak2020, boeing2025}. Parallel metaheuristics, particularly \ac{GA}, \ac{ACO}, and hybrid approaches, have demonstrated scalability and effectiveness in vehicle routing and network design problems because they can search large solution spaces efficiently, can be parallelized, and in several studies have produced strong solution quality compared with single-method baselines \citep{wens2025, alkhalifa2025, wu2024}.

Filipino commuter behavior is shaped by service convenience, walking thresholds, and strong transfer aversion \citep{upncts2021b, huang2022}. This transfer aversion is particularly pronounced, as empirical studies establish that passengers tolerate transfers only when achieving substantial time savings, otherwise prioritizing direct, uninterrupted travel time minimization \citep{harris2022, huang2022}. Another key behavioral dimension is \ac{WTW}, which refers to the maximum distance a commuter is willing to walk to or from a transit stop before choosing an alternative mode or route. \ac{WTW} thresholds ranging from 500–700 meters significantly influence stop accessibility and route choice \citep{ha2023, karesdotter2022}. These behavioral parameters must be integrated into optimization frameworks to ensure that generated routes align with actual commuter decision-making rather than theoretical ideals \citep{upncts2021b}.

Memetic algorithms (MAs) are a class of hybrid metaheuristics that embed individual local search within a population-based global search framework, combining the broad exploration of \ac{GA} with solution-level refinement to address complex combinatorial optimization problems such as vehicle routing, network design, and transit planning \citep{ong2004}. A defining characteristic of MAs is their learning mechanism: under Lamarckian learning, locally acquired improvements are not discarded but are instead directly inherited by the next generation, meaning the improved individual replaces its pre-refinement counterpart in the population \citep{ong2004}. The choice of local search operator within this Lamarckian structure is critical, as it determines the quality and direction of refinement passed forward to successive generations. \ac{ACO} is a particularly well-suited operator for this role, as it discovers high-quality paths through iterative pheromone reinforcement guided by canonical dispersal and evaporation rules, naturally biasing refinement toward high-demand, frequently traversed corridors \citep{wu2024, korzen2024}. When \ac{ACO} serves as the local search operator within a memetic \ac{GA}, its pheromone-based mutations constitute the Lamarckian step: each \ac{GA} individual undergoes \ac{ACO}-driven improvement along high-demand corridors, and the resulting solution is directly reintroduced into the population, intensifying search where it matters most \citep{yoo2024, rodriguez-esparza2024}. Parallelization of these methods across multi-core architectures reduces computational bottlenecks, enabling practical route exploration for large-scale urban networks \citep{wens2025, chaves2024}.

Building on \citet{sanchez2025}' framework, which demonstrated 15--25\% commute time reductions through parallelized agent-based route evaluation on grid-based networks, this study extends the approach by: (1) incorporating empirically calibrated passenger behavior via local survey data; (2) implementing a hybrid Memetic \ac{GA} with \ac{ACO}-inspired Lamarckian learning as the optimization engine, wherein pheromone-based local refinements are directly inherited by successive generations; (3) modeling realistic multi-route jeepney systems with capacity constraints and operational parameters; and (4) applying the framework to Iligan City's actual urban geography and demand patterns. These enhancements transition the proof-of-concept into an operationally viable, empirically grounded optimization system capable of informing real-world policy and planning decisions.}

% Chapter 1 Section 1.2: Statement of the Problem

\section{Statement of the Problem}
\noindent{Current jeepney route systems in Iligan City operate without systematic, data-driven optimization, resulting in overlapping coverage, inefficient passenger flows, and unnecessarily extended commute times. As a result, service quality is reduced and urban congestion is exacerbated. In addition, these routes are not currently evaluated against actual commuter behavior, leaving no clear means of determining whether they align with how passengers actually travel. This gap is addressed in this study through behavioral simulation and route optimization. }

% Chapter 1 Section 1.3: Research Objectives

\section{Research Objectives}
\noindent{The following are the general and specific objectives of this study.}

% Chapter 1 Section 1.3 Subsection 1.3.1: General Objective

\subsection{General Objective}
\noindent{To design, implement, and evaluate a parallelized hybrid Memetic \ac{GA} with \ac{ACO}-inspired Lamarckian learning agent-based optimization system that minimizes Total User Cost (incorporating travel, waiting, and transfer penalties, subject to vehicle capacity limits and fleet equity regularizers) across multi-route jeepney networks by integrating empirically calibrated behavioral simulation and applying the framework to Iligan City's urban transport context.}

% Chapter 1 Subsection 1.3 Subsection 1.3.2: Specific Objectives

\subsection{Specific Objectives}
\begin{enumerate}
  \item To conduct a behavioral survey among Iligan City commuters to empirically determine willingness-to-walk thresholds, transfer aversion parameters, and mode choice preferences for calibrating agent decision rules.
  \item To design and implement an agent-based passenger simulation module that models realistic commuter behaviors such as walking, waiting, riding, alighting on \ac{OSM}-derived street networks with residential and non-residential land-use distinctions based on monocentric city principles.
  \item To develop a hybrid Memetic \ac{GA} with \ac{ACO}-inspired Lamarckian learning optimization engine that evolves multi-route jeepney systems through genetic operators such as crossover, mutation, selection combined with ant colony pheromone-based local refinement to balance exploration and exploitation.
  \item To implement parallel simulation evaluation of candidate jeepney route systems across multiple processor cores, enabling scalable assessment of thousands of route configurations within feasible computation times.
  \item To benchmark the optimized jeepney system against Iligan City’s existing route network using standardized metrics such as average commute time, transfer rates, stop accessibility, route coverage and validate results through simulation-based inference calibration.
\end{enumerate}

% Chapter 1 Section 1.4: Scope and Limitations of the Research

\section{Scope and Limitations of the Research}
\noindent{This research focuses on the design, implementation, and computational evaluation of a parallelized hybrid Memetic \ac{GA} with \ac{ACO}-inspired Lamarckian learning to optimize multi-route jeepney networks. The spatial boundaries of the study are restricted to the primary road network of Iligan City, Philippines, modeled using OpenStreetMap (\ac{OSM}) data. The behavioral parameters of passenger agents—including willingness-to-walk thresholds, transfer aversion, and boarding costs—are calibrated using empirical local survey data. 

To maintain computational tractability and focus on the primary transit network design problem, the scope of this study is subject to several key limitations and exclusions:
\begin{itemize}
    \item \textbf{Static Demand Assumption:} The Origin-Destination (\ac{OD}) travel demand matrix is assumed to be static and inelastic. The study does not model dynamic demand shifts, passenger modal split adaptations (e.g., commuters shifting to private cars, tricycles, or walking as a primary mode in response to route changes), or long-term land-use changes.
    \item \textbf{Simplified Traffic and Congestion:} While historical travel times are calibrated using speed profiles retrieved from the TomTom \ac{API}, the agent-based simulation does not model dynamic micro-level traffic congestion. Physical vehicle-to-vehicle queueing, intersection delays, traffic light phases, and bottlenecks are not simulated; thus, the feedback loop between optimized jeepney routing and private vehicular traffic congestion is excluded.
    \item \textbf{User-Centric Objective Function:} The optimization engine evaluates route systems based on a multi-dimensional Total User Cost (incorporating walking, waiting, in-vehicle travel, and transfer times, with a fleet equity regularizer). It does not incorporate operator-side financial metrics, such as fuel consumption, fleet procurement costs, driver wages, fare collection efficiency, or environmental emissions.
    \item \textbf{Static Route Structures:} Evolved routes are assumed to be spatially fixed and operating on two-way streets. One-way street restrictions, dynamic route diversions, and real-time fleet dispatch/headway modifications during live operations are not considered.
    \item \textbf{Fixed Infrastructure and Safety:} The simulation assumes a hail-and-ride operational model without modeling dedicated physical stop infrastructure, boarding bays, or pedestrian safety factors. External risk factors, such as road safety, pedestrian sidewalk quality, street lighting, and traffic law enforcement, require separate modeling pipelines and are outside the scope of this work.
    \item \textbf{System-Level Evaluation:} The optimization is evaluated strictly at the system-wide network level. The performance is judged by aggregate commuter utility across the entire multi-route network rather than maximizing the profitability or service frequency of any single route in isolation.
\end{itemize}
While the modular cost formulation allows future researchers to incorporate multi-objective terms (such as operator costs or safety indicators) simply by modifying the fitness weights, the present study is deliberately bounded to user-centric travel cost minimization under a fixed fleet size of vehicles and a preset capacity of 16 passengers per vehicle.}

% Chapter 1 Section 1.5: Significance of the Research

\section{Significance of the Research}
\noindent{This research advances computer science by combining parallel hybrid metaheuristics with agent-based behavioral modeling tailored for informal public transit optimization. Theoretically, it demonstrates how Memetic \ac{GA} with \ac{ACO}-inspired Lamarckian learning algorithms and empirically calibrated agent-based simulation can be unified to address real-world, city-scale transport problems. Practically, the study provides a reproducible methodology and computational tool that can inform data-driven route planning for urban jeepney networks, integrating local commuter behavior into optimization. This contribution has potential to influence future research and applications in intelligent transportation systems, urban modeling, and scalable combinatorial optimization.}